\section{arrays}

This section documents the array and related data types. These are
dynamically sized, with the size stored at all times in the size
member. These classes are written as a C++ wrapper around a C Library
{\tt c\_arrays}. This enables the {\tt c\_arrays} library to be
implemented to make use of parallel computation and an example
implementation in C* ({\tt cs\_arrays.cs}) is provided. Unfortunately,
the function \verb|offmul| is actually the computationally dominant
routine, and this proves to be the most difficult to optimise for
parallel computers.

\subsection{array and iarray}\label{array}

\begin{verbatim}

class array
{
 public:

  int size; 
  void *list;  

  array (int s=0);
  array ( array & x);
  ~ array ();
  array  operator=( array  x); 

  operator double *();
  array  operator=( double  * x);

  /* unary minus */
  array  operator-();

  /* binary arithmetic operators */
  array   operator  +  (  array   x);
  array   operator  -  (  array   x);
  array   operator  *  (  array   x);
  array   operator  /  (  array   x);

  /* binary comparison operators */
  iarray  operator  <  (  array   x);
  iarray  operator  <=  (  array   x);
  iarray  operator  >  (  array   x);
  iarray  operator  >=  (  array   x);
  iarray  operator  ==  (  array   x);
  iarray  operator  !=  (  array   x);

  /* concatenation */
  array  operator << ( array  x);

  /* broadcast arithmetic ops - commutative equivalents defined
  elsewhere */
  array   operator  +  (  double   x);
  array   operator  -  (  double   x);
  array   operator  *  (  double   x);
  array   operator  /  (  double   x);

  /* broadcast comparison ops  - commutative equivalents defined
  elsewhere */
  iarray  operator  <  (  double   x); 
  iarray  operator  <=  (  double   x);
  iarray  operator  >  (  double   x);
  iarray  operator  >=  (  double   x);
  iarray  operator  ==  (  double   x);
  iarray  operator  !=  (  double   x);

  /* concatenate an element  - commutative equivalents defined
  elsewhere */
  array  operator << ( double  x);

  /* compound assignment */
  array   operator  += (  array   x);
  array   operator  += (  double   x); /* broadcast add to */
  array   operator  -= (  array   x);
  array   operator  -= (  double   x);
  array   operator  *= (  array   x);
  array   operator  *= (  double   x);
  array   operator  /= (  array   x);
  array   operator  /= (  double   x);
  array   operator  <<= (  array   x); /* compound concatenation */ 
  array   operator  <<= (  double   x);

  /* index operations */
#ifdef CONTIGUOUS
  double&   operator[](int);
#else
  par_addr_double operator[](int);
#endif
  par_addr_array  operator[](iarray); /* gather/scatter ops */

/* all of the previous routines have been defined for the iarray class
as well, with every occurance of array being replace by iarray, and
every occurance of double being replaced by int */

/* multiply an array by an iarray */
  inline array operator*(iarray);

/* broadcast an integer to an array */
  array operator=(int x);

/* probabilistically round an array to an iarray. See rounding below */
  operator iarray();
};  

class iarray
{
 public:
  /* all the previous methods defined for array ar also defined
     equivalently for iarray */

  /* extra iarray methods */
  /* boolean operators */
  iarray operator ! ();
  iarray operator  && (iarray x);
  iarray operator || (iarray x);
  iarray operator&=(iarray x) {return *this = *this && x;}
  iarray operator|=(iarray x) {return *this = *this || x;}

 /* probabilistic assignment - see rounding below */
  void operator=(array x);    
  void operator+=(array x);  

 /* convert to array */
  inline operator array();
}

\end{verbatim}

The implementation of the data storage is pointed to by
\verb|list|. More than one array variable may have the same
\verb|list| pointer, so that assignment does not involve copying of
the data, just copying the list pointer. This allows arrays to be
passed by value efficiently in the C++ code. In the case where {\tt
  list} points to a contiguous memory location where the data is
stored, defining {\tt CONTIGUOUS} allows obvious optimisations for
index operations.

iarray is the integer version of array. Binary operations +,*,=
etc. have been defined elementwise, otherwise a scalar value may be
broadcast over the whole array. The concatenation operation joins two
arrays to make a larger array, e.g. $\{1,2,3\} << \{4,5,6\} =
\{1,2,3,4,5,6\}$. 

\subsubsection{Indexing and Gather/Scatter}

The expression \verb|a[10]| is meant to be an {\em lvalue} (i.e. can
be assigned to) that refers to the tenth element of a. When the values
of a are internally stored in a contiguous memory location, this is
easy, with the method returning a reference. However, if the data is
stored in a distributed fashion, references are not possible, and a
function will need to be called to extract the value, or alternatively
alter the value, depending on the context. This achieved by a helper
class \verb|par_addr_double|. This can either be type cast to a
\verb|double|, in which case the function returning the value of
\verb|a[i]| is called, alternatively, it can be assigned to, in which
case the alteration function is called. If more than one array
variable refers to the same underlying data, a copy is made of the
data before an array element is altered.

A {\em gather} operation, or {\em vector index} operation is expressed
by \verb|a[i]| where {\tt i} is an {\tt iarray} variable. The
expression \verb|a[i]=x| {\em scatters} the values of {\tt x} into the
array {\tt a} according to the vector {\tt i}. Naturally, the length
of {\tt i} and {\tt x} must agree. A fresh copy of {\tt a} data is
made if needed.


\subsubsection{Rounding}

The assignment of an array variable to an iarray variable is
interesting, in that the results are rounded up or down
probabilistically. If a value 0.2 is to be assigned, then it has an
80\% chance of being rounded down to 0, and a 20\% chance of being
rounded up to 1.

\subsection{{\tt sparse\_mat}}\label{interact_array}

The {\tt sparse\_mat} data type is designed to hold the
interaction matrix \bbeta, which is generally a sparse matrix. It is
built up of array and iarray components:

\begin{verbatim}
class interact_array
{
 public:
  array diag, val;
  iarray row, col;
  sparse_mat(int s=0, int o=0)
    {diag=array(s); val=array(o); row=iarray(o); col=iarray(o);}
  array operator*(iarray& x);  /* matrix multiplication */
  sparse_mat operator=(sparse_mat x);
};
\end{verbatim}

{\tt diag} stores the diagonal components of the array, {\tt val} is
the packed list of offdiagonal values, with {\tt row} and {\tt col}
being the index lists. The only important operator defined for this
class is the matrix operation, which is defined as
\begin{verbatim}
beta*x == beta.diag*x + (beta.val*x[beta.col])[beta.row]
\end{verbatim}
but is implemented separately for efficiency reasons.

It is possible to represent the offdiagonal array differently  for
efficiency reasons. For example, if it is desired to represent the
offdiagonal elements as a dense 2D array, one can create an extra
pointer in the underlying implementation of array. Most of the time,
this pointer is NULL, but when the \verb|cs_arrays| routine
\verb|offmul| is called, it will check this pointer for the {\tt val}
array. If it is NULL, it will create the efficient
representation, otherwise it will reuse the existing one. Because the
only way the array's actual value will get out of synch with the
efficient representation is by index assignment, \verb|put_double()|
and \verb|put_double_array| (as well as obviously
\verb|delete_array()|) will need to be modified to deallocate the
efficient representation.

\subsection{Global functions}

In all of the following, if it reasonable for an iarray version to
exist, then it will.

\begin{itemize}
\item array {\tt merge}( iarray {\em mask}, array {\em a}, array {\em b} )

return an array (or iarray) {\em r} such that where {\em mask[i]==1},
{\em r[i]=a[i]}, otherwise {\em r[i]=b[i]}.

\item array {\tt pack}( array {\em x}, iarray {\em mask}, [int {\em
ntrue}])

Construct a new array from elements of  {\em x} that correspond to
where the mask is true. {\em ntrue} is an optional parameter equal to
the number of true elements of {\em mask}. If not given, then {\tt
pack} will count the true elements of {\em mask}. Give this parameter
if you are using the same {\em mask} in multiple {\tt pack} statements.

\item iarray {\tt enumerate}(iarray {\em mask})

Return the running sum of mask.

\item iarray {\tt pcoord}( int {\em size})

return [0..size-1]

\item iarray {\tt gen\_index}(iarray {\em x}) 

generate a list of species numbers, with each number appearing in the
list according to the value passed in its position. For example, if
{\em x}=\{0,0,1,2,0,1\} then the output from this program will be
\{3,4,4,6\}


\end{itemize}

\subsubsection{Reduction functions}

\begin{itemize}
\item double {\tt max}( array {\em x}, [iarray {\em mask}]) 
\item double {\tt min}( array {\em x}, [iarray {\em mask}]) 
\item double {\tt sum}( array {\em x}, [iarray {\em mask}]) 

Return the maximum, minimum and sum of an array, optionally masked.

\end{itemize}

\subsubsection{Array functions}
\begin{itemize}
\item array {\tt abs}(array {\em x})
\item iarray {\tt sign}(array {\em x})
\item array {\tt exp}(array {\em x})
\item array {\tt log}(array {\em x})
\end{itemize}

\subsubsection{Random number functions}

\begin{itemize}
\item void {\tt fillrand}(array {\em x}) 

Fill {\em x} with random numbers from the uniform distribution over $[0,1]$

\item void {\tt fillprand}(array {\em x})

Fill {\em x} with random numbers from the Poisson distribution $e^{-x}$

\item void {\tt fillgrand}(array {\em x})

Fill {\em x} with random numbers from the Gaussian distribution $e^{-x^2/4}$

\end{itemize}

\subsubsection{Stream Functions}
\begin{itemize}
\item \verb|ostream& operator<<(ostream& s, array x)|
\item \verb|ostream& operator<<(ostream& s, iarray x)|
\item \verb|ostream& operator<<(ostream& s, par_addr_double x)|
\item \verb|ostream& operator<<(ostream& s, par_addr_int x)|
\end{itemize}

